\documentclass[11pt,a4paper]{article}

\usepackage{fontspec}
\usepackage{polyglossia}

\setmainlanguage{russian}
\setotherlanguage{english}

\IfFontExistsTF{Times New Roman}
{\setmainfont{Times New Roman}}
{\setmainfont{Libertinus Serif}}

\IfFontExistsTF{Menlo}
{\newfontfamily\cyrillicfonttt{Menlo}}
{\newfontfamily\cyrillicfonttt{DejaVu Sans Mono}}

\usepackage[a4paper,
top=2cm,
bottom=2cm,
left=2.1cm,
right=2.1cm]{geometry}

\usepackage{graphicx}
\usepackage{booktabs}
\usepackage{tabularx}
\usepackage{array}
\usepackage{float}
\usepackage{caption}
\usepackage{enumitem}
\usepackage{amsmath,amssymb}
\usepackage{xurl}
\usepackage{hyperref}

\hypersetup{
    colorlinks=true,
    linkcolor=black,
    citecolor=black,
    urlcolor=blue
}

\setlength{\parindent}{1.25em}
\setlength{\parskip}{0.45em}
\setlength{\emergencystretch}{3em}

\linespread{1.04}\selectfont

\captionsetup{
    font=small,
    labelfont=bf
}

\setlength{\textfloatsep}{10pt plus 2pt minus 2pt}
\setlength{\intextsep}{10pt plus 2pt minus 2pt}

\setlist[itemize]{
    topsep=3pt,
    itemsep=2pt,
    parsep=0pt,
    leftmargin=1.6em
}

\begin{document}

\begin{titlepage}

\centering

{\large Федеральное государственное автономное образовательное учреждение высшего образования\par}

{\large НАЦИОНАЛЬНЫЙ ИССЛЕДОВАТЕЛЬСКИЙ УНИВЕРСИТЕТ\par}

{\large «ВЫСШАЯ ШКОЛА ЭКОНОМИКИ»\par}

{\large Факультет экономических наук\par}

{\large Образовательная программа «Экономика»\par}

\vfill

{\Large\bfseries
Исследование устойчивости моделей волатильности и ценообразования опционов в различных рыночных режимах
\par}

\vfill

\begin{flushright}
\textbf{Выполнила:}
\end{flushright}

\begin{flushright}
Ковалёва Мария Сергеевна
\end{flushright}

\begin{flushright}
3 курс, группа БЭК224
\end{flushright}

\vspace{1cm}

\begin{flushright}
\textbf{Научный руководитель:}\\
Сизых Наталья Васильевна\\[0.5cm]
\underline{\hspace{5cm}}
\end{flushright}

\vfill

{\large Москва, 2026 г.\par}

\end{titlepage}

\tableofcontents

\newpage

\section{Введение}

Производные финансовые инструменты возникли как ответ на потребность участников рынка в формировании более гибких и индивидуальных условий заключения сделок, а также передаче и перераспределении финансовых рисков. Самые ранние формы срочных контрактов были связаны с необходимостью заранее фиксировать условия будущей сделки, снижать неопределённость и делать торговлю более предсказуемой. В современной финансовой системе эта функция сохранилась, но инструменты стали значительно сложнее: деривативы используются для самых различных целей, таких как хеджирование, спекулятивная торговля, арбитраж, управление портфелями и оценка рыночных ожиданий \cite{Hull}.

Опционы занимают особое место среди производных инструментов, поскольку в отличие от форвардов и фьючерсов дают держателю право, а не обязанность совершить сделку с базовым активом. Такая асимметричная структура делает опционы удобным инструментом управления риском, но одновременно усложняет задачу оценки их справедливой стоимости. Цена опциона зависит не только от текущей цены базового актива, но и от страйка, срока до экспирации, процентной ставки, волатильности и стиля исполнения. Поэтому задача корректного ценообразования опционов является одной из центральных задач количественных финансов.

Практическая значимость этой задачи усиливается масштабом мирового рынка деривативов. По данным Банка международных расчётов, номинальный объём непогашенных внебиржевых деривативов на конец июня 2025 года составил около 846 трлн долларов США, а их валовая рыночная стоимость --- около 21{,}8 трлн долларов США. Такие значения показывают, что деривативы являются значимой частью современной финансовой инфраструктуры. Следовательно, вопросы оценки, устойчивости и интерпретации моделей ценообразования имеют не только теоретическое, но и прикладное значение.

\subsection*{Актуальность исследования}
\addcontentsline{toc}{subsection}{Актуальность исследования}

Актуальность данного исследования подтверждается и развитием эмпирической литературы по ценообразованию производных финансовых инструментов. В исследованиях~\cite{Bakshi}, \cite{Dumas} и~\cite{Christoffersen} подчёркивается, что сравнение моделей должно проводиться на рыночных данных, поскольку более сложная теоретическая модель не всегда автоматически даёт существенное улучшение качества оценки.

Дополнительная практическая значимость темы связана с масштабом рынка деривативов. Производные инструменты активно используются финансовыми организациями, инвесторами и участниками срочного рынка для управления рисками и оценки рыночных ожиданий. Ошибки в оценке опционов могут приводить к неверной интерпретации стоимости риска, некорректному сравнению контрактов и искажению выводов о состоянии рынка. Поэтому анализ устойчивости моделей ценообразования имеет прикладное значение даже в случае, когда рассматриваются базовые модели, используемые как бенчмарк.

В данной работе исследование проводится на опционах на фьючерс на индекс МосБиржи. Выбор индексного инструмента обусловлен тем, что Индекс МосБиржи отражает динамику наиболее ликвидных акций крупнейших российских эмитентов и выступает одним из ключевых индикаторов состояния российского рынка в целом. Именно поэтому анализ цен этих опционов позволяет рассматривать более широкий рыночный контекст. Это делает выбранный инструмент подходящим для исследования того, как качество моделей ценообразования меняется в зависимости от рыночных условий.

\subsection*{Объект исследования}
\addcontentsline{toc}{subsection}{Объект исследования}

Объектом исследования являются процессы ценообразования опционов на российском финансовом рынке.

\subsection*{Предмет исследования}
\addcontentsline{toc}{subsection}{Предмет исследования}

Предметом исследования является устойчивость моделей волатильности и ценообразования опционов в различных рыночных режимах, а также зависимость ошибок модельной оценки от срока до экспирации, денежности опциона и способа задания волатильности.

\subsection*{Проблема исследования}
\addcontentsline{toc}{subsection}{Проблема исследования}

Главная проблема данного исследования заключается в ограниченной способности моделей ценообразования опционов описывать реальные рыночные цены.

Классические модели основаны на ряде допущений и предположений о динамике базового актива, процентной ставке, волатильности и структуре рынка. В реальных рыночных данных эти предположения не выполняются: волатильность меняется во времени, доходности демонстрируют разные значения в зависимости от кластеров, присутствуют тяжёлые хвосты, качество моделей может зависеть от состояния рынка и характеристик конкретного опциона.

Более того, модели работают плохо, поскольку учитывают ограниченное количество параметров. Финансовый рынок можно рассматривать как сложную систему, в которой цена формируется под воздействием множества факторов, от количества участников рынка до макроэкономических и фундаментальных событий. Многие эти факторы являются ненаблюдаемыми, стохастическими или не поддаются достоверной оценке из-за неполноты информации. В итоге причина того, что модели ошибаются, --- переизбыток реальных переменных, а не их недостаток~\cite{Mantegna}.

Подход, который имеет смысл реализовать для решения данной проблемы, --- разработать модель моделей~\cite{Ilyinsky}. Он заключается в том, что каждый параметр общепринятой модели получается на основе работы ещё одной модели, и таким образом значение оцениваемого параметра становится наиболее приближённым к целевому. Как расширение возможно реализовать динамическую оценку: параметры будут меняться в зависимости от времени, правил пересчёта или других параметров.

\subsection*{Гипотеза исследования}
\addcontentsline{toc}{subsection}{Гипотеза исследования}

Качество аппроксимации рыночных цен опционов модельными оценками зависит от рыночного режима и ухудшается в стрессовых состояниях рынка, что требует анализа и систематизации характера и масштаба этих отклонений для дальнейшего учёта выявленных зависимостей при построении и применении моделей оценки стоимости опционов в реальных задачах.

\subsection*{Цель работы}
\addcontentsline{toc}{subsection}{Цель работы}

Цель работы --- исследовать устойчивость моделей волатильности и ценообразования опционов в различных рыночных режимах на основе сравнения модельных и рыночных цен.

\subsection*{Задачи исследования}
\addcontentsline{toc}{subsection}{Задачи исследования}

Для достижения поставленной цели в работе решаются следующие задачи:

\begin{itemize}
    \item рассмотреть теоретические основы классических моделей ценообразования опционов и подходов к оценке волатильности;

    \item описать предпосылки и ограничения классических моделей;

    \item сформировать рабочую выборку дневных наблюдений по фьючерсным опционам и соответствующему базовому активу;

    \item рассчитать дополнительные показатели для оценки цены опциона с помощью классических моделей;

    \item реализовать американскую биномиальную модель Кокса--Росса--Рубинштейна с разными способами задания волатильности;

    \item использовать европейскую модель Блэка (1976) как отдельный бенчмарк без досрочного исполнения и оценить различие между европейской и американской постановкой;

    \item сравнить модельные и рыночные цены опционов с использованием нескольких метрик качества;

    \item проанализировать ошибки моделей по режимам волатильности, срокам до экспирации и денежности опционов;

    \item провести формальные проверки ошибок: тесты средней ошибки, парные сравнения функций потерь и описательную регрессию абсолютной ошибки на характеристики режима и контракта;

    \item проверить устойчивость результатов на отдельных подвыборках и при изменении процентной ставки;

    \item сформулировать выводы о применимости рассмотренных моделей и ограничениях простой постановки с одной оценкой волатильности на контракт.
\end{itemize}

\section{Анализ современных методов и моделей оценки волатильности и ценообразования опционов}

\subsection{Классические модели}

В настоящей работе рассматриваются три классические модели ценообразования опционов. Биномиальная модель Кокса--Росса--Рубинштейна~\cite{CRR} --- численный метод, который строит дискретное дерево цен базового актива и допускает досрочное исполнение, что необходимо для американских опционов. Модель Блэка--Шоулза~\cite{BlackScholes,Merton} --- аналитическая формула для европейских опционов, более простая в вычислении, но не позволяющая учесть право раннего исполнения. Модель Блэка 1976 года~\cite{Black} --- адаптация Блэка--Шоулза для опционов на фьючерсы и форварды, используемая в данной работе как европейский аналитический эталон.

Все три модели принимают волатильность как входной параметр. Именно волатильность в наибольшей степени определяет цену опциона при зафиксированных прочих параметрах, и именно её оценка является наиболее сложной задачей на практике. В данной работе волатильность не задаётся константой, а оценивается отдельно --- с помощью исторических и GARCH-моделей. Таким образом, реализуется подход «модели моделей», описанный в разделе «Проблема исследования»: каждый параметр ценообразующей модели получается на выходе отдельной модели-оценщика.

Для оценки волатильности используются модели класса ARCH/GARCH~\cite{Engle,Bollerslev}. Модель ARCH (Engle, 1982) описывает условную дисперсию как функцию прошлых квадратов остатков, улавливая кластеризацию волатильности --- явление, при котором периоды высокой волатильности сменяются аналогичными периодами. Модель GARCH (Bollerslev, 1986) расширяет ARCH, добавляя зависимость условной дисперсии от её собственных прошлых значений, что позволяет описывать более длительную «память» волатильности с меньшим числом параметров. В данной работе используется спецификация GARCH(1,1) как наиболее распространённая и устойчивая на практике.

\subsection{Биномиальная модель Кокса--Росса--Рубинштейна}

Биномиальная модель Кокса--Росса--Рубинштейна представляет собой численный метод ценообразования опционов, предложенный в 1979 году. В её основе лежит построение дискретного биномиального дерева возможных путей цены базового актива. Модель является конечно-разностным приближением к геометрическому броуновскому движению и при неограниченном увеличении числа шагов сходится к аналитической формуле Блэка--Шоулза.

\subsubsection{Построение биномиального дерева}

Срок до экспирации $T$ разбивается на $N$ равных шагов длиной $\Delta t = T/N$. На каждом шаге цена базового актива умножается либо на коэффициент роста $u$, либо на коэффициент снижения $d$:

\begin{equation}
u = \exp(\sigma\sqrt{\Delta t}), \qquad d = \frac{1}{u},
\end{equation}

где $\sigma$ --- волатильность базового актива. Параметризация обеспечивает симметрию $u \cdot d = 1$ и согласована с дисперсией лог-доходностей геометрического броуновского движения.

\subsubsection{Риск-нейтральная вероятность}

Риск-нейтральная вероятность восходящего движения определяется из условия отсутствия арбитража. Для базового актива с безрисковой ставкой $r$:

\begin{equation}
q = \frac{e^{r\Delta t} - d}{u - d}.
\end{equation}

Для опциона на фьючерс условие мартингальности фьючерсной цены приводит к упрощённому выражению:

\begin{equation}
q = \frac{1 - d}{u - d}.
\end{equation}

\subsubsection{Ценообразование и досрочное исполнение}

На конечном узле дерева стоимость опциона равна его внутренней стоимости: для опциона на покупку $H = \max(F - K, 0)$, для опциона на продажу $H = \max(K - F, 0)$.

Далее применяется обратная индукция от экспирации к текущему моменту. Для европейского опциона:

\begin{equation}
V = e^{-r\Delta t}\!\left(q V^{\uparrow} + (1-q) V^{\downarrow}\right).
\end{equation}

Для американского опциона в каждом узле дополнительно проверяется выгодность немедленного исполнения:

\begin{equation}
V = \max\!\left( e^{-r\Delta t}\!\left(q V^{\uparrow} + (1-q) V^{\downarrow}\right),\; H(F, K) \right).
\end{equation}

Именно это свойство делает биномиальную модель предпочтительной для ценообразования американских опционов: аналитические формулы, в частности Блэка--Шоулза и Блэка 1976 года, не позволяют учесть возможность досрочного исполнения.

\subsection{Модель Блэка--Шоулза}

Модель Блэка--Шоулза (Блэка--Шоулза--Мертона) является одной из базовых моделей ценообразования опционов и представляет собой фундаментальный результат современной финансовой теории. Она задаёт аналитическую формулу для оценки стоимости европейских опционов на основе ряда предположений о динамике базового актива и свойствах финансового рынка.

\subsubsection{Динамика базового актива}

В основе модели лежит предположение о том, что цена базового актива $S_t$ следует геометрическому броуновскому движению:

\begin{equation}
dS_t = \mu S_t dt + \sigma S_t dW_t,
\end{equation}

где:

\begin{itemize}
    \item $\mu$ --- ожидаемая доходность актива;
    \item $\sigma$ --- волатильность;
    \item $W_t$ --- стандартное броуновское движение.
\end{itemize}

Из данного предположения следует, что логарифм цены актива имеет нормальное распределение, а сама цена --- логнормальное распределение.

\subsubsection{Основные предположения модели}

Модель Блэка--Шоулза основана на следующих ключевых предположениях:

\begin{itemize}
    \item отсутствуют арбитражные возможности;
    
    \item рынок является совершенным, то есть отсутствуют транзакционные издержки и налоги;
    
    \item торговля осуществляется непрерывно;
    
    \item возможно заимствование и инвестирование под постоянную безрисковую ставку $r$;
    
    \item волатильность $\sigma$ постоянна во времени;
    
    \item базовый актив не выплачивает дивиденды (в базовой версии модели);
    
    \item динамика цены актива описывается геометрическим броуновским движением.
\end{itemize}

Эти предположения позволяют получить аналитическое решение задачи ценообразования опциона.

\subsubsection{Идея модели: репликация и отсутствие арбитража}

Ключевая идея модели заключается в построении безрискового портфеля, состоящего из опциона и базового актива. За счёт динамического хеджирования устраняется стохастическая составляющая доходности портфеля.

Доходность такого портфеля должна совпадать с безрисковой ставкой процента:

\begin{equation}
\Pi = V - \Delta S,
\end{equation}

где:
\begin{itemize}
    \item $V$ --- цена опциона;
    \item $\Delta = \frac{\partial V}{\partial S}$ --- дельта опциона.
\end{itemize}

Применение леммы Ито и условия отсутствия арбитража приводит к дифференциальному уравнению Блэка--Шоулза:

\begin{equation}
\frac{\partial V}{\partial t}
+
\frac{1}{2}\sigma^2 S^2 \frac{\partial^2 V}{\partial S^2}
+
rS \frac{\partial V}{\partial S}
-
rV
=
0.
\end{equation}

\subsubsection{Формула Блэка--Шоулза}

Решение данного уравнения для европейского опциона на покупку имеет вид:

\begin{equation}
C = S_0 N(d_1) - K e^{-rT} N(d_2),
\end{equation}

где:

\begin{equation}
d_1 =
\frac{
\ln\left(\frac{S_0}{K}\right)
+
\left(r + \frac{\sigma^2}{2}\right)T
}{
\sigma \sqrt{T}
},
\end{equation}

\begin{equation}
d_2 = d_1 - \sigma \sqrt{T},
\end{equation}

а $N(\cdot)$ --- функция распределения стандартного нормального закона.

Для европейского опциона на продажу используется формула:

\begin{equation}
P = K e^{-rT} N(-d_2) - S_0 N(-d_1).
\end{equation}

\subsubsection{Интерпретация формулы}

Формула Блэка--Шоулза имеет следующую экономическую интерпретацию:

\begin{itemize}
    \item $S_0 N(d_1)$ отражает приведённую стоимость ожидаемой цены базового актива;
    
    \item $K e^{-rT} N(d_2)$ отражает приведённую стоимость выплаты по страйку;
    
    \item разность этих величин определяет стоимость опциона как ожидаемую выгоду от исполнения контракта.
\end{itemize}

Также модель может быть интерпретирована в рамках риск-нейтрального подхода, где цена опциона равна дисконтированному математическому ожиданию его выплаты:

\begin{equation}
C = e^{-rT}\mathbb{E}^Q[(S_T - K)^+].
\end{equation}

\subsubsection{Входные параметры модели}

Модель принимает на вход следующие параметры:

\begin{itemize}
    \item $S_0$ --- текущая цена базового актива;
    
    \item $K$ --- страйк опциона;
    
    \item $T$ --- время до экспирации;
    
    \item $r$ --- безрисковая процентная ставка;
    
    \item $\sigma$ --- волатильность.
\end{itemize}

\subsubsection{Выход модели}

Результатом работы модели являются:

\begin{itemize}
    \item цена опциона;
    
    \item чувствительности опциона (показатели чувствительности): $\Delta$, $\Gamma$, $\Theta$, $\nu$, $\rho$.
\end{itemize}

\subsubsection{Роль волатильности}

Волатильность является ключевым параметром модели Блэка--Шоулза. Именно она определяет степень неопределённости будущей динамики базового актива и существенно влияет на стоимость опциона.

На практике вместо исторической волатильности часто используется подразумеваемая волатильность, извлекаемая из наблюдаемых рыночных цен опционов.

Таким образом, модель Блэка--Шоулза часто применяется в обратном направлении: не для вычисления цены опциона, а для определения подразумеваемой волатильности, которая отражает рыночные ожидания относительно будущей волатильности.

\subsubsection{Ограничения модели Блэка--Шоулза}

Несмотря на фундаментальную роль модели Блэка--Шоулза, её предположения существенно упрощают реальную динамику финансовых рынков. В результате модель демонстрирует систематические отклонения от наблюдаемых рыночных цен.

\paragraph{Постоянство волатильности}

Одним из ключевых предположений модели является постоянство волатильности $\sigma$. Однако эмпирические данные показывают, что волатильность:

\begin{itemize}
    \item изменяется во времени;
    
    \item демонстрирует кластеризацию;
    
    \item зависит от состояния рынка.
\end{itemize}

Это приводит к тому, что одна фиксированная волатильность не может корректно описать цены опционов на разных страйках и сроках.

\paragraph{Улыбка волатильности}

В реальных рыночных данных наблюдается зависимость подразумеваемой волатильности от страйка и срока до экспирации:

\begin{equation}
\sigma_{\text{подр}} = f(K, T),
\end{equation}

в то время как модель Блэка--Шоулза предполагает:

\begin{equation}
\sigma = \text{пост.}
\end{equation}

Эмпирически это проявляется в виде улыбки волатильности или перекоса волатильности, что означает систематическую ошибку модели при оценке опционов вне денег.

\paragraph{Нормальность лог-доходностей}

Модель предполагает нормальность логарифмических доходностей. Однако на практике наблюдаются:

\begin{itemize}
    \item тяжёлые хвосты;
    
    \item асимметрия распределения;
    
    \item повышенная вероятность экстремальных событий.
\end{itemize}

Это означает, что модель недооценивает риск сильных рыночных движений.

\paragraph{Отсутствие скачков}

В модели предполагается непрерывная динамика цены:

\begin{equation}
dS_t = \mu S_t dt + \sigma S_t dW_t,
\end{equation}

что исключает резкие скачки цен. В реальности финансовые рынки подвержены:

\begin{itemize}
    \item новостным шокам;
    
    \item гэпам;
    
    \item кризисным движениям.
\end{itemize}

Это приводит к недооценке стоимости опционов, чувствительных к экстремальным событиям.

\paragraph{Постоянство процентной ставки}

Модель предполагает фиксированную безрисковую ставку $r$, тогда как в реальности процентные ставки изменяются во времени и обладают собственной временной структурой.

\paragraph{Совершенность рынка}

Предполагается отсутствие:

\begin{itemize}
    \item транзакционных издержек;
    
    \item ограничений на короткие продажи;
    
    \item ликвидностных ограничений.
\end{itemize}

Кроме того, модель требует непрерывного хеджирования, что невозможно реализовать на практике.

\paragraph{Последствия для практического применения}

В результате перечисленных ограничений:

\begin{itemize}
    \item модель систематически ошибается в оценке опционов;
    
    \item ошибка зависит от страйка, срока и рыночного режима;
    
    \item качество модели снижается в стрессовых рыночных условиях.
\end{itemize}

Это делает необходимым использование более сложных моделей, а также мотивирует эмпирическую проверку устойчивости модели в различных рыночных режимах.





\subsection{Модель Блэка 1976 года}

Модель Блэка 1976 года является адаптацией формулы Блэка--Шоулза для ценообразования европейских опционов на фьючерсы и форвардные контракты. В отличие от базовой модели, где стохастической переменной служит спотовая цена актива $S_t$, здесь в качестве базового актива используется фьючерсная цена $F_t$.

\subsubsection{Динамика фьючерсной цены}

Поскольку фьючерсный контракт не требует первоначальных инвестиций, в риск-нейтральной мере ожидаемая доходность фьючерса равна нулю. Модель предполагает, что фьючерсная цена следует геометрическому броуновскому движению без дрейфа:

\begin{equation}
dF_t = \sigma F_t\, dW_t.
\end{equation}

\subsubsection{Формула Блэка 1976 года}

Для европейского опциона на покупку:

\begin{equation}
C = e^{-rT}\!\left[F_0 N(d_1) - K N(d_2)\right],
\end{equation}

для европейского опциона на продажу:

\begin{equation}
P = e^{-rT}\!\left[K N(-d_2) - F_0 N(-d_1)\right],
\end{equation}

где:

\begin{equation}
d_1 = \frac{\ln\!\left(\dfrac{F_0}{K}\right) + \dfrac{\sigma^2}{2}T}{\sigma\sqrt{T}},
\qquad
d_2 = d_1 - \sigma\sqrt{T},
\end{equation}

а $N(\cdot)$ --- функция распределения стандартного нормального закона.

\subsubsection{Входные параметры}

Модель принимает на вход следующие параметры:

\begin{itemize}
    \item $F_0$ --- текущая фьючерсная цена;
    \item $K$ --- страйк опциона;
    \item $T$ --- время до экспирации;
    \item $r$ --- безрисковая процентная ставка;
    \item $\sigma$ --- волатильность фьючерсной цены.
\end{itemize}

\subsubsection{Ограничения модели}

Модель Блэка 1976 года наследует ключевые предположения формулы Блэка--Шоулза: постоянство волатильности, нормальность логарифмических доходностей фьючерсной цены и непрерывность торговли. Принципиальным ограничением в контексте настоящего исследования является то, что модель рассчитана исключительно на европейские опционы и не учитывает возможность досрочного исполнения. Для опционов на фьючерсы с американским стилем исполнения необходимо использование численных методов, в частности биномиальной модели.

\subsection{Альтернативные модели ценообразования}

Рассмотренные выше модели составляют классическое ядро теории ценообразования опционов и служат базовым эталоном для сравнения. Вместе с тем литература предлагает ряд более сложных подходов, направленных на устранение ключевых ограничений классических моделей.

Среди наиболее известных расширений --- модели со стохастической волатильностью (Heston, 1993 и SABR), модели со скачками (Merton, 1976 и Kou, 2002), а также модели с локальной волатильностью (Dupire, 1994). Эти подходы позволяют воспроизводить улыбку волатильности, тяжёлые хвосты распределения и другие стилизованные факты, которые не описываются классическими моделями при постоянной волатильности. Каждое из расширений, однако, требует дополнительных данных для калибровки, усложняет реализацию и интерпретацию результатов.

Детальное рассмотрение этих моделей выходит за рамки настоящей работы. Исследование сосредоточено на биномиальной модели Кокса--Росса--Рубинштейна с тремя спецификациями волатильности и модели Блэка 1976 года как аналитическом эталоне без досрочного исполнения.

\subsection{Волатильность в ценообразовании опционов}

Волатильность является одним из ключевых параметров моделей ценообразования опционов. Она характеризует степень изменчивости цены базового актива и напрямую влияет на стоимость опциона. Чем выше ожидаемая волатильность, тем выше вероятность значительных ценовых движений, а следовательно --- тем выше стоимость опциона.

В рамках моделей типа Блэка--Шоулза волатильность выступает в качестве параметра, определяющего распределение возможных будущих значений цены базового актива. На практике корректная оценка волатильности является одной из центральных задач количественных финансов.

\subsubsection{Историческая волатильность}

Историческая волатильность представляет собой оценку волатильности, рассчитанную на основе исторических данных о доходностях базового актива.

Как правило, историческая волатильность рассчитывается как выборочное стандартное отклонение логарифмических доходностей с приведением к годовому масштабу. Для окна длиной $n$ торговых дней:

\begin{equation}
r_t = \ln \left( \frac{S_t}{S_{t-1}} \right),
\end{equation}

\begin{equation}
\hat{\sigma}^2 = \frac{1}{n-1}\sum_{j=0}^{n-1}\!\left(r_{t-j} - \bar{r}\right)^2,
\qquad
\bar{r} = \frac{1}{n}\sum_{j=0}^{n-1} r_{t-j},
\end{equation}

\begin{equation}
HV_n = \sqrt{252}\cdot\hat{\sigma}.
\end{equation}

Множитель $\sqrt{252}$ приводит дневную волатильность к годовой при допущении о 252 торговых днях в году.

На практике историческая волатильность может рассчитываться на различных временных окнах, например 21 или 63 торговых дня. Различные окна позволяют учитывать краткосрочную или более долгосрочную динамику рынка.

Основным преимуществом исторической волатильности является простота расчёта. Однако данный подход использует исключительно прошлые данные и не учитывает рыночные ожидания относительно будущей волатильности.

\subsubsection{Реализованная волатильность}

Реализованная волатильность характеризует фактически реализованную волатильность за некоторый период времени. В отличие от исторической волатильности, которая используется как оценка будущей волатильности, реализованная волатильность представляет собой последующую характеристику уже произошедшей динамики рынка.

В современных исследованиях реализованная волатильность часто используется как эталон при анализе качества прогнозов волатильности и моделей прогнозирования волатильности.

\subsubsection{Подразумеваемая волатильность}

Подразумеваемая волатильность представляет собой значение волатильности, при котором модельная цена опциона совпадает с наблюдаемой рыночной ценой.

Формально подразумеваемая волатильность определяется как решение уравнения:

\begin{equation}
C^{\text{рын}} = C^{\text{мод}}(\sigma_{\text{подр}}),
\end{equation}

где:
\begin{itemize}
    \item $C^{\text{рын}}$ --- наблюдаемая рыночная цена опциона;
    
    \item $C^{\text{мод}}$ --- модельная цена опциона;
    
    \item $\sigma_{\text{подр}}$ --- подразумеваемая волатильность.
\end{itemize}

Уравнение не имеет аналитического решения, поэтому подразумеваемая волатильность находится численно, например методом Ньютона--Рафсона:

\begin{equation}
\sigma^{(k+1)} = \sigma^{(k)} - \frac{C^{\text{мод}}(\sigma^{(k)}) - C^{\text{рын}}}{\mathcal{V}(\sigma^{(k)})},
\end{equation}

где $\mathcal{V} = \partial C^{\text{мод}}/\partial\sigma$ --- вега опциона, то есть чувствительность модельной цены к волатильности.

Таким образом, подразумеваемая волатильность позволяет выразить рыночную цену опциона через единый параметр волатильности. На практике она широко используется как универсальный язык рынка опционов и позволяет сравнивать опционы с различными страйками и сроками до экспирации.

В отличие от исторической волатильности, подразумеваемая волатильность отражает ожидания участников рынка относительно будущей неопределённости и потенциальных ценовых движений.

\subsubsection{Улыбка волатильности и перекос волатильности}

Одним из наиболее известных эмпирических эффектов на рынке опционов является улыбка волатильности --- зависимость подразумеваемой волатильности от страйка опциона.

Модель Блэка--Шоулза предполагает постоянную волатильность:

\begin{equation}
\sigma = \text{пост.},
\end{equation}

однако в реальных рыночных данных подразумеваемая волатильность зависит как от страйка, так и от срока до экспирации:

\begin{equation}
\sigma_{\text{подр}} = f(K, T).
\end{equation}

Для многих рынков характерна ситуация, при которой опционы далеко вне денег обладают более высокой подразумеваемой волатильностью, чем опционы около денег. На фондовых рынках часто наблюдается асимметричная структура перекоса волатильности, связанная с повышенным спросом на защитные опционы на продажу.

Наличие улыбки волатильности и перекоса волатильности свидетельствует о том, что предположения модели Блэка--Шоулза не полностью соответствуют реальной динамике финансовых рынков.

\subsubsection{Временная структура волатильности}

Подразумеваемая волатильность зависит не только от страйка, но и от срока до экспирации опциона. Зависимость подразумеваемой волатильности от времени до погашения называется временной структурой волатильности.

Структура волатильности может принимать различные формы в зависимости от состояния рынка и ожиданий участников. В периоды стрессовых состояний краткосрочная подразумеваемая волатильность часто существенно возрастает, отражая повышенную неопределённость на коротком горизонте.

Анализ временной структуры волатильности играет важную роль при построении поверхностей волатильности и моделировании цен опционов.

\subsection{Модели GARCH}

Одним из ключевых ограничений модели Блэка--Шоулза является предположение о постоянстве волатильности. Эмпирические данные показывают, что волатильность финансовых активов изменяется во времени и обладает свойством кластеризации волатильности: периоды высокой волатильности часто следуют за периодами высокой волатильности, а спокойные периоды --- за спокойными.

Для моделирования подобных эффектов используются модели условной гетероскедастичности, наиболее известной из которых является модель GARCH.

\subsubsection{Кластеризация волатильности}

Одной из характерных особенностей финансовых временных рядов является наличие кластеризации волатильности. Несмотря на то, что сами доходности часто обладают слабой автокорреляцией, величины абсолютных или квадратов доходностей демонстрируют устойчивую зависимость во времени.

Это означает, что большие ценовые движения имеют тенденцию группироваться во времени, формируя периоды повышенной рыночной турбулентности.

\subsubsection{Модель GARCH(1,1)}

Модель GARCH, то есть обобщённая авторегрессионная условная гетероскедастичность описывает условную дисперсию доходностей как функцию прошлых шоков и прошлых значений волатильности.

Базовая модель GARCH(1,1) имеет вид:

\begin{equation}
r_t = \mu + \varepsilon_t,
\end{equation}

\begin{equation}
\varepsilon_t = \sigma_t z_t,
\end{equation}

\begin{equation}
\sigma_t^2 =
\omega
+
\alpha \varepsilon_{t-1}^2
+
\beta \sigma_{t-1}^2,
\end{equation}

где:
\begin{itemize}
    \item $\sigma_t^2$ --- условная дисперсия;
    
    \item $\omega$ --- долгосрочный уровень волатильности;
    
    \item $\alpha$ --- чувствительность к новым рыночным шокам;
    
    \item $\beta$ --- влияние прошлой волатильности;
    
    \item $z_t$ --- случайная величина со стандартным нормальным распределением.
\end{itemize}

Для стационарности процесса необходимо выполнение условия:

\begin{equation}
\alpha + \beta < 1.
\end{equation}

Параметр $(\alpha+\beta)$ называется персистентностью: при значениях близких к единице волатильность возвращается к долгосрочному уровню медленно. Безусловная (долгосрочная) дисперсия определяется как:

\begin{equation}
\bar{\sigma}^2 = \frac{\omega}{1 - \alpha - \beta}.
\end{equation}

\subsubsection{Прогнозирование волатильности}

Прогноз условной дисперсии на $h$ шагов вперёд имеет аналитическую форму:

\begin{equation}
\hat{\sigma}^2_{t+h|t} = \bar{\sigma}^2 + (\alpha + \beta)^{h-1}\!\left(\sigma^2_t - \bar{\sigma}^2\right).
\end{equation}

При $h=1$ это однодневный прогноз $\hat{\sigma}^2_{t+1|t} = \omega + \alpha\varepsilon_t^2 + \beta\sigma_t^2$. Для использования в моделях ценообразования однодневный прогноз приводится к годовому масштабу:

\begin{equation}
\hat{\sigma}_{\text{ann},t} = \sqrt{252}\,\hat{\sigma}_{t+1|t}.
\end{equation}

\subsubsection{Использование модели GARCH в задачах ценообразования опционов}

Модель GARCH не используется напрямую для расчёта цены опциона. Вместо этого она применяется для прогнозирования будущей волатильности базового актива, которая затем подставляется как параметр $\sigma$ в модель ценообразования. В рамках данного исследования прогноз $\hat{\sigma}_{\text{ann},t}$ подставляется в биномиальную модель Кокса--Росса--Рубинштейна. Это позволяет сравнить качество прогнозной условной волатильности с историческими оценками и исследовать их влияние на точность аппроксимации рыночных цен.






\section{Методология исследования}

\subsection{Постановка исследовательской задачи}

Эмпирическая часть работы направлена на проверку того, насколько модели ценообразования способны воспроизводить наблюдаемые расчётные цены фьючерсных опционов в разных состояниях рынка. Рыночная цена рассматривается как наблюдаемое значение, сформированное участниками торгов, а модельная цена --- как теоретическая оценка, полученная при заданных предпосылках о динамике базового фьючерса, процентной ставке, волатильности и праве раннего исполнения.

Такой подход согласуется с логикой гипотезы эффективного рынка, согласно которой рыночные цены агрегируют доступную участникам информацию и их ожидания относительно будущей динамики базового актива и риска. Поэтому наблюдаемая биржевая расчётная цена используется в работе как эмпирический бенчмарк для сравнения моделей. При этом она не интерпретируется как абсолютно истинная фундаментальная стоимость опциона, поскольку на неё могут влиять ликвидность, особенности биржевого расчёта, дисбаланс спроса и предложения и временные рыночные отклонения.

Ключевое методологическое уточнение состоит в том, что рассматриваемые контракты являются опционами на фьючерсы и допускают американский стиль исполнения. Поэтому основное сравнение проводится между американскими биномиальными моделями Кокса--Росса--Рубинштейна с разными способами задания волатильности. Европейская модель Блэка 1976 года используется отдельно как эталонная нижняя оценка без раннего исполнения, а не как полноценная конкурирующая модель для основного ранжирования.

\subsection{Подготовка данных и рабочей выборки}

В качестве исходной базы используется очищенный набор дневных наблюдений по опционам серии ММ на фьючерс MXI за период с 3 января 2024 года по 19 мая 2026 года. Фьючерс MXI (Mini-MosExchange Index) --- это фьючерсный контракт на Индекс МосБиржи (IMOEX), торгуемый на Срочном рынке Московской биржи. Серия ММ --- это серия опционов на данный фьючерс с ежемесячными экспирациями. В датасете содержится 19 574 строки, 92 уникальных опционных инструмента и 5 базовых фьючерсных контрактов. Даты экспирации находятся в диапазоне с 16 апреля 2026 года по 17 июня 2027 года.

Рыночная цена опциона задаётся расчётной ценой биржи. В качестве цены базового актива используется расчётная цена соответствующего фьючерса, а не значение фондового индекса. Это важно, поскольку для фьючерсных опционов модель должна строиться относительно фьючерсной цены. Для каждого наблюдения рассчитываются срок до экспирации в календарных днях, время до экспирации в годах, денежность как отношение цены фьючерса к страйку и логарифмическая денежность.

В рабочую выборку включаются только наблюдения с корректными датами, положительной рыночной ценой, положительной ценой фьючерса, положительным страйком, положительным сроком до экспирации, распознанным типом опциона и непустым идентификатором базового фьючерса. После фильтрации остаётся 19 565 наблюдений, то есть 99,954\% исходной базы. Общая сопоставимая выборка для сравнения трёх американских моделей содержит 18 058 наблюдений: именно для них одновременно доступны все три оценки волатильности.

\subsection{Оценка волатильности}

Для каждого базового фьючерса строится временной ряд дневных логарифмических доходностей:

\begin{equation}
r_t = \ln\left(\frac{F_t}{F_{t-1}}\right),
\end{equation}

где $F_t$ --- расчётная цена фьючерса в день $t$. На этой основе рассчитываются три входных показателя волатильности.

Первый показатель --- историческая волатильность за 21 торговый день. Она рассчитывается как скользящее стандартное отклонение дневных доходностей с приведением к годовому масштабу:

\begin{equation}
\sigma_{21,t} = \sqrt{252}\,\operatorname{sd}(r_{t-20},\ldots,r_t),
\end{equation}

где $\operatorname{sd}(\cdot)$ --- выборочное стандартное отклонение. Экономически эта оценка отражает краткосрочную динамику рынка и быстрее реагирует на недавние изменения цен. Её слабая сторона состоит в повышенной зашумлённости, особенно для опционов с дальними экспирациями.

Второй показатель --- историческая волатильность за 63 торговых дня:

\begin{equation}
\sigma_{63,t} = \sqrt{252}\,\operatorname{sd}(r_{t-62},\ldots,r_t).
\end{equation}

Эта оценка соответствует более длинному горизонту и является более устойчивой. Она менее чувствительна к краткосрочным выбросам, что особенно важно для выборки, где преобладают опционы с большим сроком до экспирации.

Третий показатель --- прогнозная условная волатильность по модели GARCH(1,1). Для каждого базового фьючерса оценивается модель с нулевым условным средним и нормальным распределением стандартизованных шоков:

\begin{equation}
r_t = \varepsilon_t, \qquad \varepsilon_t = \sigma_t z_t,
\end{equation}

\begin{equation}
\sigma_t^2 = \omega + \alpha \varepsilon_{t-1}^2 + \beta \sigma_{t-1}^2.
\end{equation}

Модель оценивается на расширяющемся окне после накопления не менее 63 наблюдений. Параметры переоцениваются через 21 торговый день, а между переоценками условная дисперсия обновляется рекурсивно с использованием новых доходностей. Полученный однодневный прогноз дисперсии переводится в годовую волатильность умножением на $\sqrt{252}$.

\subsection{Рыночные режимы и группировка наблюдений}

Режимы волатильности строятся по распределению 21-дневной исторической волатильности базовых фьючерсов. Наблюдения делятся на три группы: низкая, средняя и высокая волатильность. Границы задаются первым и вторым терцилями распределения. Дополнительно рассчитывается 21-дневная доходность базового фьючерса, и по её знаку формируется направление рынка. Этот признак используется как вспомогательная характеристика режима.

Для анализа структуры ошибок также используются группы по сроку до экспирации и денежности. По сроку выделяются интервалы 0--7 дней, 8--30 дней, 31--90 дней и более 91 дня. По денежности выделяются глубоко вне денег, вне денег, около денег, в деньгах и глубоко в деньгах. Такая группировка позволяет проверить, связаны ли ошибки не только с общим режимом рынка, но и с положением контракта на опционной поверхности.

\subsection{Модели ценообразования}

\subsubsection{Американская биномиальная модель Кокса--Росса--Рубинштейна}

Основной моделью ценообразования в данном исследовании является биномиальная модель Кокса--Росса--Рубинштейна для американского опциона на фьючерс. Выбор этой модели обусловлен тем, что рассматриваемые контракты допускают досрочное исполнение, которое не учитывается в аналитических формулах. В расчётах используется 100 шагов дерева. Процентная ставка в базовой спецификации фиксируется на уровне 16\% годовых.

В работе оцениваются три спецификации, различающиеся способом задания входной волатильности:

\begin{itemize}
    \item модель Кокса--Росса--Рубинштейна с 21-дневной исторической волатильностью;
    \item модель Кокса--Росса--Рубинштейна с 63-дневной исторической волатильностью;
    \item модель Кокса--Росса--Рубинштейна с прогнозной волатильностью GARCH.
\end{itemize}

\subsubsection{Европейская модель Блэка 1976 года}

В качестве дополнительного ориентира рассчитывается европейская модель Блэка 1976 года с 63-дневной исторической волатильностью. Поскольку эта модель не учитывает досрочное исполнение, она не включается в основное ранжирование американских спецификаций, а используется для оценки премии раннего исполнения --- разности между ценой КРР и ценой европейской формулы при одинаковой волатильности:

\begin{equation}
\Pi^{R} = V^{\text{КРР}}_{63} - V^{\text{Блэк}}_{63}.
\end{equation}

\subsection{Расчёт ошибок и критерии качества}

Для каждого наблюдения рассчитывается знаковая ошибка:

\begin{equation}
e_i = \widehat{P}_i - P_i,
\end{equation}

где $\widehat{P}_i$ --- модельная цена, а $P_i$ --- рыночная расчётная цена. Положительная ошибка означает переоценку опциона моделью, отрицательная --- недооценку.

Далее строятся абсолютная, квадратичная и относительная ошибки. В итоговых таблицах используются четыре стандартные метрики:

\begin{align}
\text{MAE}  &= \frac{1}{n}\sum_{i=1}^{n}|e_i|, &
\text{RMSE} &= \sqrt{\frac{1}{n}\sum_{i=1}^{n}e_i^{2}}, \\[4pt]
\overline{e} &= \frac{1}{n}\sum_{i=1}^{n}e_i, &
\text{MAPE} &= \frac{1}{n}\sum_{i=1}^{n}\left|\frac{e_i}{P_i}\right|.
\end{align}

Средняя абсолютная ошибка (MAE) показывает средний масштаб отклонения в пунктах цены, RMSE сильнее штрафует крупные промахи, средняя знаковая ошибка $\overline{e}$ позволяет выявить систематическое смещение, а MAPE нормирует ошибку на полную рыночную цену.

Все основные сравнения проводятся на общей выборке, где доступны цены всех трёх американских спецификаций. Это исключает ситуацию, при которой различия в качестве объясняются не моделью, а разным составом наблюдений.

В дополнение к стандартным метрикам вводится авторский показатель --- нормированная ошибка по временно́й стоимости (TVNAE, Time-Value Normalized Absolute Error). Для каждого наблюдения он определяется как:

\begin{equation}
\text{TVNAE}_i = \frac{|\hat{P}_i - P_i|}{TV_i},
\qquad
TV_i = P_i - H(F_i, K_i),
\end{equation}

где $TV_i$ --- временна́я стоимость опциона (time value), а $H(F, K)$ --- его внутренняя стоимость: $\max(F - K,\, 0)$ для колла и $\max(K - F,\, 0)$ для пута.

Введение TVNAE обусловлено особенностью данной выборки. Стандартные абсолютные метрики (MAE, RMSE) измеряют отклонение в пунктах цены. Однако значительная часть наблюдений --- это длинные опционы глубоко в деньгах. Их цена состоит преимущественно из внутренней стоимости, которая определяется разностью фьючерсной цены и страйка и одинаково воспроизводится всеми моделями независимо от точности прогноза волатильности. В результате большая абсолютная ошибка на таком контракте отражает не столько качество волатильностной модели, сколько масштаб контракта. Ошибка в 100 пунктов на опционе с временно́й стоимостью 80 пунктов означает, что модель ошиблась на 125\% от той части цены, за которую она «отвечает». Та же ошибка при временно́й стоимости 500 пунктов --- это отклонение в 20\%.

TVNAE нормирует ошибку именно на временну́ю стоимость --- компоненту, чувствительную к волатильности, сроку до экспирации и процентной ставке. Это даёт более честное сравнение моделей по их ключевой функции. Стандартная относительная ошибка MAPE нормирует на полную цену, а не на опциональную компоненту: для контрактов глубоко в деньгах, где временна́я стоимость составляет, например, 5\% от общей цены, MAPE и MAE дают схожую информацию, в то время как TVNAE корректно выделяет компоненту, которую модель действительно оценивает.

Расчёт TVNAE ограничивается наблюдениями с $TV_i \geq 10$ пунктов, что исключает нестабильность при малых значениях знаменателя. После применения порога в анализ включается 16\,961 наблюдение из 18\,058 в общей сопоставимой выборке (93{,}9\%).

\subsection{Эконометрическая проверка ошибок}

Формальная проверка проводится в три этапа. Сначала для каждой американской модели оценивается средняя ошибка и проверяется гипотеза о её равенстве нулю. Спецификация теста:

\begin{equation}
e_{i,t} = \mu + \varepsilon_{i,t},
\end{equation}

где $e_{i,t}$ --- ошибка $i$-го контракта в день $t$, $\mu$ --- проверяемый параметр. Стандартные ошибки кластеризуются по торговой дате $t$, что учитывает внутридневную зависимость ошибок, обусловленную общими рыночными условиями.

Затем выполняются парные тесты Диболда--Мариано. Для каждой пары моделей $A$ и $B$ строится дневная средняя разность функций потерь:

\begin{equation}
d_t = \bar{L}^{A}_{t} - \bar{L}^{B}_{t},
\end{equation}

где $\bar{L}^{A}_{t}$ --- среднее значение функции потерь (абсолютной или квадратичной) по всем контрактам дня $t$ для модели $A$. Нулевая гипотеза $\mathbb{E}[d_t] = 0$ проверяется регрессией:

\begin{equation}
d_t = \delta + u_t,
\end{equation}

со стандартными ошибками, устойчивыми к гетероскедастичности и автокорреляции (HAC с пятью лагами по Newey--West). Это позволяет учесть серийную зависимость дневных разностей.

На третьем этапе оценивается описательная OLS-регрессия абсолютной ошибки на характеристики наблюдения:

\begin{equation}
|e_i| = \beta_0
        + \beta_1\,\mathbf{1}[\text{high\_vol}]_i
        + \beta_2\,\mathrm{DTE}_i
        + \beta_3\,\bigl|\ln(F_i/K_i)\bigr|
        + \beta_4\,\mathbf{1}[\text{put}]_i
        + \varepsilon_i,
\end{equation}

где $\mathbf{1}[\text{high\_vol}]_i$ --- индикатор режима высокой волатильности, $\mathrm{DTE}_i$ --- срок до экспирации в календарных днях, $|\ln(F_i/K_i)|$ --- абсолютная логарифмическая денежность. Стандартные ошибки кластеризуются по торговой дате. Эта регрессия не является моделью цены --- её цель описать, какие характеристики связаны с величиной ошибки.

Дополнительно проводится проверка чувствительности к процентной ставке: базовая ставка 16\% сравнивается со значениями 12\% и 20\%. Также выполняются проверки устойчивости на подвыборках: без крайних хвостов денежности, только для краткосрочных контрактов, только для длинных контрактов и без режима высокой волатильности.

\section{Эмпирическое исследование устойчивости моделей}

\subsection{Описание данных}

Исходная база содержит 19 574 наблюдения. В ней представлены 13 217 опционов на покупку (колл) и 6 357 опционов на продажу (пут). Наиболее многочисленным базовым контрактом является фьючерс MXI с экспирацией в июне 2026 года: на него приходится 11 709 строк. Остальные наблюдения распределены между более дальними фьючерсами с экспирациями в декабре 2026 года, сентябре 2026 года, марте 2027 года и июне 2027 года.

Критически важные для моделирования поля заполнены полностью: дата торгов, дата экспирации, тип опциона, страйк, расчётная цена опциона, идентификатор базового фьючерса и расчётная цена фьючерса. Основные пропуски сосредоточены в торговых полях открытия, максимума, минимума, закрытия и объёма, поэтому они не используются в модели ценообразования. Такой выбор делает исследование основанным на расчётных ценах, а не на внутридневных сделках.

Средний срок до экспирации в полной базе равен примерно 480 календарным дням, медианный срок --- 475 дней. Это означает, что выборка существенно смещена в сторону длинных опционов. Средняя денежность, определяемая как отношение цены фьючерса к страйку, равна 1,2045, а медианная --- 1,1623. Следовательно, значительная часть наблюдений находится в зоне опционов в деньгах и глубоко в деньгах.

\subsection{Сравнение американских моделей}

Основное сравнение проводится между тремя американскими спецификациями на общей выборке из 18 058 наблюдений. Результаты приведены в таблице \ref{tab:american-models}.

\begin{table}[H]
\centering
\caption{Ошибки американских моделей на общей выборке}
\label{tab:american-models}
\small
\begin{tabular}{lrrrrrr}
\toprule
Модель & $N$ & \begin{tabular}{@{}c@{}}MAE,\\пункты\end{tabular} & \begin{tabular}{@{}c@{}}RMSE,\\пункты\end{tabular} & \begin{tabular}{@{}c@{}}Средняя\\ошибка\end{tabular} & \begin{tabular}{@{}c@{}}Медианная\\абс. ошибка\end{tabular} & \begin{tabular}{@{}c@{}}TVNAE\\(средн.)\end{tabular} \\
\midrule
КРР, волатильность GARCH   & 18 058 & 107,01 & 145,76 & $-$78,14 & 82,88 & \textbf{0{,}660} \\
КРР, волатильность 63 дня  & 18 058 & 107,15 & 140,63 & $-$89,72 & 84,45 & 0{,}696 \\
КРР, волатильность 21 день & 18 058 & 119,31 & 156,41 & $-$95,26 & 96,84 & 0{,}763 \\
\bottomrule
\end{tabular}
\end{table}

По стандартным метрикам (MAE, RMSE) лидирует модель с 63-дневной исторической волатильностью по RMSE (140,63) и GARCH по MAE (107,01), причём разница между ними по MAE составляет лишь 0,14 пункта. Модель с 21-дневной волатильностью уступает обеим: MAE~=~119,31, RMSE~=~156,41. Все три модели систематически недооценивают рыночные цены (отрицательная средняя ошибка); наименьшее смещение у GARCH ($-$78,14).

Авторская метрика TVNAE даёт принципиально иную картину. Разрыв между GARCH (0{,}660) и 63-дневной волатильностью (0{,}696) составляет 3{,}6 процентного пункта --- против 0{,}1 пункта по MAE. Это означает, что стандартные метрики маскировали преимущество GARCH: они усредняют ошибку по всем контрактам, включая длинные глубоко-в-деньгах опционы, у которых большая часть цены --- внутренняя стоимость, одинаково хорошо воспроизводимая всеми моделями независимо от точности прогноза волатильности. TVNAE нормирует ошибку только на ту компоненту цены (временну́ю стоимость), за которую модель «отвечает», и за счёт этого обнаруживает, что GARCH точнее оценивает опциональную составляющую.

\textit{Вывод.} Авторская метрика TVNAE позволяет делать более содержательные выводы, чем стандартные MAE и RMSE, в условиях выборки, смещённой в сторону длинных и глубоко-в-деньгах опционов. Именно TVNAE корректно выявляет лидерство GARCH-спецификации: преимущество скрыто в MAE за счёт доминирования контрактов, где волатильностная компонента мала относительно полной цены.

\subsection{Европейский эталон и премия раннего исполнения}

\begin{table}[H]
\centering
\caption{Сравнение американских моделей и европейского эталона}
\label{tab:euro-benchmark}
\small
\begin{tabular}{lrrrrr}
\toprule
Модель & $N$ & \begin{tabular}{@{}c@{}}MAE,\\пункты\end{tabular} & \begin{tabular}{@{}c@{}}RMSE,\\пункты\end{tabular} & \begin{tabular}{@{}c@{}}Средняя\\ошибка\end{tabular} & \begin{tabular}{@{}c@{}}Медиана\\абс. ошибки\end{tabular} \\
\midrule
КРР, 63-дн. HV  & 18 058 & 107,15 & 140,63 & $-$89,72  & 84,45  \\
КРР, GARCH      & 18 058 & 107,01 & 145,76 & $-$78,14  & 82,88  \\
КРР, 21-дн. HV  & 18 058 & 119,31 & 156,41 & $-$95,26  & 96,84  \\
\midrule
Блэк-76, 63-дн. HV & 18 058 & 156,16 & 211,64 & $-$149,00 & 109,97 \\
\bottomrule
\end{tabular}
\end{table}

\begin{table}[H]
\centering
\caption{Средняя премия раннего исполнения КРР(63-дн.) $-$ Блэк-76(63-дн.) по режиму и сроку}
\label{tab:early-exercise}
\small
\begin{tabular}{lrrrr}
\toprule
Режим волатильности & 0--7 дн. & 8--30 дн. & 31--90 дн. & 91+ дн. \\
\midrule
Низкая  & 0,05 & 0,22 &  2,13 &  40,66 \\
Средняя & 0,05 & 0,80 &  2,81 &  64,83 \\
Высокая & ---  & ---  &  ---  & 105,94 \\
\bottomrule
\end{tabular}
\end{table}

При добавлении европейской модели Блэка 1976 года с 63-дневной волатильностью MAE возрастает до 156,16 пункта, а RMSE --- до 211,64. Средняя ошибка составляет $-$149,00, что в 1,7 раза больше по модулю, чем у лучшей американской спецификации.

Этот результат не следует интерпретировать как провал по качеству волатильностной оценки. Главная причина --- отсутствие досрочного исполнения. Это подтверждает таблица~\ref{tab:early-exercise}: для контрактов со сроком до 30 дней премия раннего исполнения пренебрежимо мала (0,05--0,80 пункта), однако для длинных контрактов (91+ дней) она резко возрастает --- от 40,66 пункта в режиме низкой волатильности до 105,94 пункта в режиме высокой. Кроме того, 24,5\% цен Блэк-76 в выборке оказались ниже внутренней стоимости соответствующего контракта, что само по себе является признаком нарушения нижней границы опциональной стоимости. Среднее значение премии раннего исполнения по всей выборке --- 59,28 пункта (медиана --- 9,87). Это подтверждает необходимость использовать американскую модель как основную для данных контрактов.

Экономическая природа этой премии состоит в следующем. Держатель американского опциона вправе исполнить контракт в любой момент и немедленно получить внутреннюю стоимость, тогда как держатель европейского вынужден ждать до экспирации. Для пута в деньгах досрочное исполнение позволяет высвободить капитал, равный страйку минус фьючерсная цена, и немедленно инвестировать эту сумму по безрисковой ставке --- тем самым временна́я стоимость денег напрямую конвертируется в ценность права раннего исполнения. Чем длиннее оставшийся срок и чем выше процентная ставка, тем дороже эта возможность. Именно это объясняет нелинейное нарастание премии в таблице~\ref{tab:early-exercise}: при сроке 91+ дней и ставке 16\% выгода от немедленного получения кэша на горизонте свыше квартала становится весомой, и европейская формула, не способная её учесть, систематически занижает теоретическую цену.

\subsection{Ошибки по режимам волатильности}

\begin{table}[H]
\centering
\caption{MAE американских моделей по режимам волатильности}
\label{tab:mae-regime}
\small
\begin{tabular}{lrrr}
\toprule
Модель & Низкая волатильность & Средняя волатильность & Высокая волатильность \\
\midrule
КРР, GARCH    & \textbf{77,49} & \textbf{118,20} & 141,00 \\
КРР, 63-дн.   & 82,08          & 128,26          & \textbf{122,83} \\
КРР, 21-дн.   & 104,12         & 141,81          & 117,76 \\
\bottomrule
\end{tabular}
\end{table}

В режимах низкой и средней волатильности лидирует GARCH-спецификация (MAE~=~77,49 и 118,20 соответственно). В режиме высокой волатильности картина меняется: GARCH ухудшается до 141,00 и уступает обеим историческим оценкам, тогда как 63-дневная (122,83) и 21-дневная (117,76) волатильность показывают меньшие значения MAE. Следовательно, прогнозная условная волатильность полезна в спокойных и умеренных режимах, но менее устойчива в наиболее напряжённом сегменте.

\subsection{Ошибки по сроку до экспирации и денежности}

\begin{table}[H]
\centering
\caption{MAE по сроку до экспирации}
\label{tab:mae-dte}
\small
\begin{tabular}{lrrrr}
\toprule
Модель & 0--7 дн. & 8--30 дн. & 31--90 дн. & 91+ дн. \\
\midrule
КРР, GARCH  & \textbf{1,88} & \textbf{4,87} & \textbf{24,64} & 122,53 \\
КРР, 63-дн. & 3,53          & 9,60          & 35,39          & \textbf{121,04} \\
КРР, 21-дн. & 2,77          & 8,79          & 35,94          & 135,28 \\
\bottomrule
\end{tabular}
\end{table}

\begin{table}[H]
\centering
\caption{MAE по денежности}
\label{tab:mae-moneyness}
\small
\begin{tabular}{lrrrrr}
\toprule
Модель & Глубоко OTM & OTM & ATM & ITM & Глубоко ITM \\
\midrule
КРР, GARCH  & \textbf{47,40}  & \textbf{67,81}  & \textbf{68,72}  & \textbf{109,11} & 125,56 \\
КРР, 63-дн. & 61,52           & 90,66           & 88,55           & 115,21          & \textbf{111,48} \\
КРР, 21-дн. & 63,38           & 97,05           & 94,12           & 129,23          & 125,30 \\
\bottomrule
\end{tabular}
\end{table}

По срокам наблюдается сильная нелинейность. На коротких контрактах (0--7 дней) ошибки всех моделей малы (1,88--3,53 пункта), при этом GARCH значительно точнее на горизонтах 8--30 и 31--90 дней. Основная масса ошибок сосредоточена в сегменте 91+ дней: здесь GARCH (122,53) и 63-дневная историческая волатильность (121,04) практически равны, тогда как 21-дневная (135,28) заметно хуже. По денежности: GARCH лидирует во всех зонах, кроме глубоко-в-деньгах, где 63-дневная историческая оценка (111,48) незначительно лучше. Крупные ошибки у глубоко-в-деньгах контрактов обусловлены не столько волатильностным прогнозом, сколько доминированием внутренней стоимости --- эту компоненту все модели описывают одинаково.

\subsection{Нормированная ошибка по временно́й стоимости (TVNAE)}

Расчёт TVNAE проводится на 16\,961 наблюдении (93{,}9\% общей выборки) после исключения контрактов с временно́й стоимостью ниже 10 пунктов. Результаты существенно уточняют выводы, полученные на основе абсолютных метрик.

На общей выборке GARCH-спецификация показывает заметное преимущество: среднее TVNAE равно 0{,}660 против 0{,}696 у 63-дневной исторической волатильности и 0{,}763 у 21-дневной. Разрыв между GARCH и 63-дневной волатильностью составляет 5{,}2 процентного пункта по TVNAE --- против практически нулевого преимущества (0{,}1 пункта) по MAE. Таким образом, GARCH лучше оценивает именно временну́ю стоимость опционов: это преимущество было скрыто в абсолютных метриках за счёт доминирования глубоко в деньгах контрактов с большой внутренней стоимостью.

\begin{table}[H]
\centering
\caption{TVNAE: общая выборка ($TV \geq 10$ пунктов)}
\label{tab:tvnae-overall}
\small
\begin{tabular}{lrrrr}
\toprule
Модель & $N$ & \begin{tabular}{@{}c@{}}Среднее\\TVNAE\end{tabular} & \begin{tabular}{@{}c@{}}Медианное\\TVNAE\end{tabular} & \begin{tabular}{@{}c@{}}Среднее\\знак. TVE\end{tabular} \\
\midrule
КРР, волатильность GARCH   & 16\,961 & 0{,}660 & 0{,}624 & $-$0{,}464 \\
КРР, волатильность 63 дня  & 16\,961 & 0{,}696 & 0{,}681 & $-$0{,}532 \\
КРР, волатильность 21 день & 16\,961 & 0{,}763 & 0{,}778 & $-$0{,}566 \\
\bottomrule
\end{tabular}
\end{table}

Разрез по волатильностным режимам выявляет принципиальную асимметрию GARCH-модели. В режимах низкой и средней волатильности GARCH уверенно лидирует: 0{,}629 против 0{,}712 у 63-дневной и 0{,}840 у 21-дневной в низком режиме. В среднем режиме --- 0{,}650 против 0{,}689 и 0{,}738 соответственно. Однако в режиме высокой волатильности картина меняется: GARCH ухудшается до 0{,}716, уступая 63-дневной (0{,}680) и 21-дневной (0{,}680) волатильности. Это подтверждает и конкретизирует ранее сделанный вывод: прогнозная условная волатильность добавляет ценность в спокойных и умеренных условиях, но в стрессовых режимах теряет это преимущество.

\begin{table}[H]
\centering
\caption{TVNAE по режимам волатильности}
\label{tab:tvnae-vol}
\small
\begin{tabular}{lrrr}
\toprule
Модель & Низкая волатильность & Средняя волатильность & Высокая волатильность \\
\midrule
КРР, волатильность GARCH   & 0{,}629 & 0{,}650 & 0{,}716 \\
КРР, волатильность 63 дня  & 0{,}712 & 0{,}689 & 0{,}680 \\
КРР, волатильность 21 день & 0{,}840 & 0{,}738 & 0{,}680 \\
\bottomrule
\end{tabular}
\end{table}

Разрез по сроку до экспирации наиболее ярко демонстрирует преимущество GARCH. На контрактах со сроком до 7 дней GARCH-спецификация достигает TVNAE 0{,}189 --- против 0{,}594 у 63-дневной и 0{,}401 у 21-дневной. На горизонте 8--30 дней --- 0{,}288 против 0{,}653 и 0{,}550 соответственно. Таким образом, GARCH воспроизводит временну́ю стоимость краткосрочных опционов в 2--3 раза точнее, чем исторические аналоги. Для интервала 31--90 дней преимущество GARCH сохраняется: 0{,}527 против 0{,}776 и 0{,}777. На длинных контрактах (более 91 дня) все три спецификации сближаются: 0{,}682 (GARCH), 0{,}688 (63-дн.) и 0{,}766 (21-дн.).

\begin{table}[H]
\centering
\caption{TVNAE по сроку до экспирации}
\label{tab:tvnae-dte}
\small
\begin{tabular}{lrrrr}
\toprule
Модель & 0--7 дн. & 8--30 дн. & 31--90 дн. & 91+ дн. \\
\midrule
КРР, волатильность GARCH   & 0{,}189 & 0{,}288 & 0{,}527 & 0{,}682 \\
КРР, волатильность 63 дня  & 0{,}594 & 0{,}653 & 0{,}776 & 0{,}688 \\
КРР, волатильность 21 день & 0{,}401 & 0{,}550 & 0{,}777 & 0{,}766 \\
\bottomrule
\end{tabular}
\end{table}

TVNAE уточняет общий вывод об эквивалентности GARCH и 63-дневной исторической волатильности, полученный на основе MAE. Эквивалентность наблюдается только на длинных контрактах, которые количественно доминируют в данной выборке. На всех остальных сегментах --- коротких, среднесрочных, ATM, OTM --- GARCH воспроизводит временну́ю стоимость существенно точнее. Практический смысл этого разграничения состоит в следующем: если портфель сконцентрирован в длинных контрактах, оба подхода дают сопоставимый результат. Если же в портфеле присутствуют короткие или около-денег опционы, GARCH предпочтительнее.

\subsection{Формальные тесты ошибок}

\begin{table}[H]
\centering
\caption{Тест средней ошибки (константа регрессии $e_i$ на 1, SE кластеризованы по дате)}
\label{tab:me-test}
\small
\begin{tabular}{lrrrrr}
\toprule
Модель & $N$ & ME & SE & $t$ & $p$ \\
\midrule
КРР, GARCH   & 18\,058 & $-$78,14 & 4,73 & $-$16,52 & $<$0,001 \\
КРР, 63-дн.  & 18\,058 & $-$89,72 & 3,62 & $-$24,77 & $<$0,001 \\
КРР, 21-дн.  & 18\,058 & $-$95,26 & 4,11 & $-$23,16 & $<$0,001 \\
\bottomrule
\end{tabular}
\end{table}

\begin{table}[H]
\centering
\caption{Парные тесты Диболда--Мариано (дневные средние разности потерь, HAC-5)}
\label{tab:dm-test}
\small
\begin{tabular}{lllrr}
\toprule
Сравнение & Функция потерь & $N_{\text{дней}}$ & $t$ & $p$ \\
\midrule
21-дн. vs 63-дн. & MAE & 489 & 3,88 & 0,0003 \\
21-дн. vs 63-дн. & MSE & 489 & 3,86 & 0,0003 \\
21-дн. vs GARCH  & MAE & 489 & 1,43 & 0,153  \\
21-дн. vs GARCH  & MSE & 489 & 1,30 & 0,195  \\
63-дн. vs GARCH  & MAE & 489 & 0,30 & 0,770  \\
63-дн. vs GARCH  & MSE & 489 & 0,90 & 0,370  \\
\bottomrule
\end{tabular}
\end{table}

\begin{table}[H]
\centering
\caption{OLS регрессия $|e_i|$ на характеристики наблюдения (КРР, 63-дн. HV; SE кластеризованы по дате)}
\label{tab:ols-reg}
\small
\begin{tabular}{lrrrr}
\toprule
Переменная & Коэф. & SE & $t$ & $p$ \\
\midrule
\multicolumn{5}{l}{\textit{Спецификация 1: один регрессор}} \\
Константа                   & 82,08 & 7,30 &  11,24 & $<$0,001 \\
$\mathbf{1}[\text{high\_vol}]$ & 21,45 & 7,73 &   2,77 & 0,006    \\
$R^2$ & \multicolumn{4}{l}{0,019} \\
\midrule
\multicolumn{5}{l}{\textit{Спецификация 2: добавлены DTE и денежность}} \\
Константа                   & 160,37 & 17,82 &  9,00 & $<$0,001 \\
$\mathbf{1}[\text{high\_vol}]$ & $-$17,60 & 6,24 & $-$2,82 & 0,005 \\
DTE (дни)                   &    0,181 & 0,030 &   5,96 & $<$0,001 \\
$|\ln(F/K)|$                & $-$206,66 & 28,29 & $-$7,30 & $<$0,001 \\
$R^2$ & \multicolumn{4}{l}{0,248} \\
\midrule
\multicolumn{5}{l}{\textit{Спецификация 3: добавлен тип опциона}} \\
Константа                   & 175,39 & 17,51 &  10,02 & $<$0,001 \\
$\mathbf{1}[\text{high\_vol}]$ & $-$15,06 & 5,87 & $-$2,57 & 0,011 \\
DTE (дни)                   &    0,162 & 0,030 &   5,37 & $<$0,001 \\
$|\ln(F/K)|$                & $-$189,19 & 28,06 & $-$6,74 & $<$0,001 \\
$\mathbf{1}[\text{put}]$    &  $-$28,29 & 5,56  & $-$5,09 & $<$0,001 \\
$R^2$ & \multicolumn{4}{l}{0,267} \\
\bottomrule
\end{tabular}
\end{table}

Все три американские спецификации статистически значимо недооценивают рыночные цены: гипотеза о нулевой средней ошибке отвергается во всех случаях (таблица~\ref{tab:me-test}). Наименьшее по модулю смещение --- у GARCH ($-$78,14).

Парные тесты (таблица~\ref{tab:dm-test}) показывают, что 21-дневная волатильность статистически значимо хуже 63-дневной ($p = 0{,}0003$). Разница между 63-дневной и GARCH статистически незначима ни по MAE ($p = 0{,}770$), ни по MSE ($p = 0{,}370$): на общей выборке оба подхода статистически эквивалентны.

Описательная регрессия (таблица~\ref{tab:ols-reg}) выявляет структуру ошибки. В однофакторной спецификации индикатор высокой волатильности положителен и значим. Однако после добавления DTE и денежности (спецификация 2) его знак меняется на отрицательный: при прочих равных высокий режим не является самостоятельным источником роста ошибки. Эффект режима в простом разрезе объясняется составом наблюдений --- в высоком режиме преобладают длинные контракты. Главный структурный предиктор ошибки --- срок до экспирации: каждый дополнительный день добавляет около 0,18 пункта к абсолютной ошибке.

\subsection{Проверки устойчивости}

\begin{table}[H]
\centering
\caption{MAE по подвыборкам (проверки устойчивости)}
\label{tab:robustness}
\small
\begin{tabular}{lrrrr}
\toprule
Подвыборка & $N$ & КРР, GARCH & КРР, 63-дн. & КРР, 21-дн. \\
\midrule
Полная выборка             & 18\,058 & \textbf{107,01} & 107,15 & 119,31 \\
Без глубоких хвостов       & 17\,221 & \textbf{87,59}  & 92,76  & 102,47 \\
Только короткие (DTE $<$ 91) &  2\,077 & \textbf{22,27}  & 29,92  & 29,41  \\
Только длинные (DTE $\geq$ 91) & 15\,981 & 122,71 & \textbf{121,04} & 135,28 \\
Без режима высокой волатильности & 14\,267 & \textbf{87,59}  & 92,76  & 102,47 \\
\bottomrule
\end{tabular}
\end{table}

\begin{table}[H]
\centering
\caption{Чувствительность к ставке (КРР, 63-дн. HV)}
\label{tab:rate-sensitivity}
\small
\begin{tabular}{lrrrr}
\toprule
Ставка & MAE & RMSE & ME & Медиана AE \\
\midrule
12\% & 104,57 & 138,21 & $-$85,40 & 81,93 \\
16\% (базовая) & 107,15 & 140,63 & $-$89,72 & 84,45 \\
20\% & 109,84 & 143,11 & $-$93,82 & 87,12 \\
\bottomrule
\end{tabular}
\end{table}

Таблица~\ref{tab:robustness} демонстрирует, что основной вывод не зависит от конкретной подвыборки. GARCH лидирует по MAE на полной выборке, без глубоких хвостов и на коротких контрактах. На длинных контрактах (DTE $\geq$ 91) 63-дневная волатильность незначительно лучше (121,04 против 122,71), однако разница находится в пределах статистической погрешности.

Чувствительность к процентной ставке (таблица~\ref{tab:rate-sensitivity}) умеренная: снижение ставки с 16\% до 12\% уменьшает MAE на 2,58 пункта. Качественные выводы о ранжировании моделей при всех трёх уровнях ставки сохраняются.

\section{Заключение}

\subsection{Ответ на гипотезу и проблему исследования}

Все три американские спецификации систематически недооценивают рыночные расчётные цены. Ни одна из моделей с единственной оценкой волатильности не воспроизводит рыночные цены без систематической погрешности: плоская поверхность волатильности является принципиальным источником смещения.

Гипотеза исследования --- о том, что качество аппроксимации зависит от рыночного режима и ухудшается в периоды повышенной волатильности --- не отвергается в описательном разрезе, однако многофакторный анализ существенно уточняет её содержание. Рыночный режим сам по себе не является самостоятельной причиной роста ошибок: его простой описательный эффект объясняется тем, что в периоды высокой волатильности в выборке преобладают длинные контракты. Главный структурный предиктор ошибки --- срок до экспирации. Один параметр волатильности не способен описать ни временну́ю структуру цен, ни улыбку волатильности, что и порождает нарастающее смещение на длинных горизонтах.

\subsection{Сравнение подходов к оценке волатильности}

На полной выборке качество GARCH и 63-дневной исторической волатильности статистически неразличимо. Они лидируют по разным критериям: GARCH лучше по типичной ошибке, 63-дневная --- устойчивее к крупным промахам. Разграничение становится очевидным при разрезе по сроку: на коротких и среднесрочных контрактах GARCH существенно точнее, поскольку его прогноз условной дисперсии улавливает актуальный уровень рыночной турбулентности. На длинных контрактах преимущество GARCH исчезает --- прогноз на один день вперёд плохо отражает волатильность на горизонте нескольких сотен дней. В режиме высокой волатильности GARCH также ухудшается относительно конкурентов, что согласуется с его известным свойством переоценивать персистентность текущих шоков.

21-дневная историческая волатильность статистически значимо хуже обоих альтернатив. Для данной выборки с преобладанием длинных контрактов короткое окно наблюдений слишком шумное и не отражает горизонт, на котором реализуется стоимость опциона.

\subsection{Роль досрочного исполнения: американская модель против европейской}

Выбор типа модели оказывается важнее выбора оценки волатильности. Европейская формула Блэка 1976 года при одинаковой волатильности даёт существенно б\'{о}льшую ошибку, чем американская биномиальная модель, --- разрыв между ними многократно превышает разницу между любыми двумя волатильностными спецификациями. Это иерархический вывод: неправильный выбор класса модели наносит больший ущерб, чем неоптимальная оценка волатильности.

Премия досрочного исполнения нелинейно нарастает с ростом срока и режима волатильности. На коротких контрактах она пренебрежимо мала. На длинных --- экономически существенна и усиливается в стрессовых условиях, когда вероятность выгодного досрочного исполнения наиболее высока.

\subsection{Систематическое смещение и его природа}

Отрицательная средняя ошибка всех трёх моделей указывает на факторы, выходящие за рамки постановки с единой волатильностью. Первая причина --- плоская поверхность волатильности: модель с единственным значением $\sigma$ не воспроизводит улыбку подразумеваемой волатильности, характерную для реального рынка. Вторая --- расчётные цены биржи могут содержать ликвидностную премию. Третья --- фиксированная ставка не учитывает срочную структуру, и использование более точной оценки ставки ослабляет недооценку, хотя и не устраняет её полностью.

\subsection{Ограничения исследования}

В качестве рыночной цены использована расчётная цена биржи, а не котировка или цена последней сделки --- на неликвидных инструментах эти значения могут расходиться. Выборка смещена в сторону длинных опционов в деньгах, что обусловливает специфический профиль ошибок и ограничивает перенесение выводов на рынки с иной структурой. Единая оценка волатильности без поверхности является принципиальным архитектурным ограничением: устранить систематическое смещение в рамках данной постановки невозможно без перехода к более гибким моделям.

\subsection{Рекомендации по применению и направления развития}

\textbf{Выбор типа модели.} Для американских фьючерсных опционов необходимо использовать американскую биномиальную модель. Применение европейской формулы систематически занижает цену длинных и глубоко в деньгах контрактов, причём эффект усиливается в периоды высокой волатильности.

\textbf{Выбор оценки волатильности.} 63-дневная историческая волатильность является надёжной базовой оценкой для широкого спектра контрактов. GARCH целесообразно применять для краткосрочных контрактов и в периоды низкой и средней волатильности. 21-дневная историческая волатильность нецелесообразна для портфелей с преобладанием длинных опционов.

\textbf{Учёт процентной ставки.} Использование срочной структуры ставок вместо единого константного значения позволяет снизить ошибку и ослабить систематическую недооценку.

\textbf{Направления развития.} Для снижения систематического смещения необходим переход к моделям с поверхностью подразумеваемой волатильности. Перспективными направлениями также являются учёт ликвидностных факторов и применение моделей со стохастической волатильностью или скачками --- классов, теоретически мотивированных наблюдаемой кластеризацией волатильности и тяжёлыми хвостами распределения доходностей.

\begin{thebibliography}{99}

\bibitem{Hull}
Hull J. C. Options, Futures, and Other Derivatives. --- 10th ed. --- Pearson, 2018.

\bibitem{Black}
Black F. The Pricing of Commodity Contracts // Journal of Financial Economics. --- 1976. --- Vol. 3. --- P. 167--179.

\bibitem{BlackScholes}
Black F., Scholes M. The Pricing of Options and Corporate Liabilities // Journal of Political Economy. --- 1973. --- Vol. 81, No. 3. --- P. 637--654.

\bibitem{Merton}
Merton R. Theory of Rational Option Pricing // Bell Journal of Economics and Management Science. --- 1973. --- Vol. 4, No. 1. --- P. 141--183.

\bibitem{Bollerslev}
Bollerslev T. Generalized Autoregressive Conditional Heteroskedasticity // Journal of Econometrics. --- 1986. --- Vol. 31. --- P. 307--327.

\bibitem{CRR}
Cox J. C., Ross S. A., Rubinstein M. Option Pricing: A Simplified Approach // Journal of Financial Economics. --- 1979. --- Vol. 7, No. 3. --- P. 229--263.

\bibitem{Engle}
Engle R. F. Autoregressive Conditional Heteroscedasticity with Estimates of the Variance of United Kingdom Inflation // Econometrica. --- 1982. --- Vol. 50, No. 4. --- P. 987--1007.

\bibitem{Hamilton}
Hamilton J. D. A New Approach to the Economic Analysis of Nonstationary Time Series and the Business Cycle // Econometrica. --- 1989. --- Vol. 57, No. 2. --- P. 357--384.

\bibitem{Mantegna}
Mantegna R. N., Stanley H. E. Introduction to Econophysics: Correlations and Complexity in Finance. --- Cambridge University Press, 1999.

\bibitem{Ilyinsky}
Ильинский К. В зеркале супермоделей. --- М.: Фазис, 2003.

\bibitem{Bakshi}
Bakshi G., Cao C., Chen Z. Empirical Performance of Alternative Option Pricing Models // Journal of Finance. --- 1997. --- Vol. 52, No. 5. --- P. 2003--2049.

\bibitem{Dumas}
Dumas B., Fleming J., Whaley R. Implied Volatility Functions: Empirical Tests // Journal of Finance. --- 1998. --- Vol. 53, No. 6. --- P. 2059--2106.

\bibitem{Christoffersen}
Christoffersen P., Jacobs K. The Importance of the Loss Function in Option Valuation // Journal of Financial Economics. --- 2004. --- Vol. 72. --- P. 291--318.

\bibitem{Cont}
Cont R. Empirical Properties of Asset Returns: Stylized Facts and Statistical Issues // Quantitative Finance. --- 2001. --- Vol. 1, No. 2. --- P. 223--236.

\bibitem{Tsay}
Tsay R. Analysis of Financial Time Series. --- 3rd ed. --- Wiley, 2010.

\bibitem{Shreve}
Shreve S. Stochastic Calculus for Finance II: Continuous-Time Models. --- Springer, 2004.

\bibitem{Haug}
Haug E. G. The Complete Guide to Option Pricing Formulas. --- McGraw-Hill, 2007.

\bibitem{Weithers}
Weithers T. Quantitative Finance. --- Wiley, 2007.

\end{thebibliography}

\end{document}
